% ---
% Dedicatória
% ---
\begin{dedicatoria}
   \vspace*{\fill}
   \centering
   \noindent
   \textit{Este trabalho é dedicado a minha família e amigos\\
   que sempre me apoiaram nos momentos mais difíceis.} \vspace*{\fill}
\end{dedicatoria}
% ---

% ---
% Agradecimentos
% ---
\begin{agradecimentos}
Agradeço a minha esposa, Maria Clara, por todo apoio e pela base que me deu durante toda a minha caminhada. À minha mãe, agradeço por ter me mostrado a importância dos estudos e por ser um pilar em quem eu me tornei. 

Agradeço em especial ao meu orientador José Antonio Moreira Xexéo, pela paciência e pela orientação ao longo de toda a pesquisa, Agradeço também a Gabriela Moutinho de Souza Dias, Anderson F. Pereira dos Santos e Ronaldo Ribeiro Goldschmidt pelo apoio a este trabalho. 
Agradeço também ao corpo docente do Instituto Militar de Engenharia pelos ensinamentos que recebi ao longo da minha formação.


\end{agradecimentos}
% ---

% ---
% Epígrafe
% ---
\begin{epigrafe}
    \vspace*{\fill}
	\begin{flushright}
		\textit{``Não vos amoldeis às estruturas deste mundo, \\
		mas transformai-vos pela renovação da mente, \\
		a fim de distinguir qual é a vontade de Deus: \\
		o que é bom, o que Lhe é agradável, o que é perfeito.''\\
		(Bíblia Sagrada, Romanos 12, 2)}
	\end{flushright}
\end{epigrafe}
% ---

% ---
% RESUMOS
% ---

% resumo em português
\setlength{\absparsep}{18pt} % ajusta o espaçamento dos parágrafos do resumo
\begin{resumo}
\SingleSpacing
%iniciomudança
A padronização da família Ascon pelo National Institute of Standards and Technology (NIST), formalizada na publicação NIST SP 800-232, consolidou a Criptografia Leve (Lightweight Cryptography - LWC) como resposta à demanda por cifragem autenticada em dispositivos com restrições de processamento, memória e energia. No entanto, a literatura de identificação de algoritmos criptográficos por aprendizado de máquina reporta acurácias elevadas sob protocolos que permitem vazamento de informação entre treinamento e teste, como o reaproveitamento de chaves e a seleção de características fora do laço de validação, e não avalia esquemas AEAD leves padronizados ou finalistas de padronização sob cenário ciphertext-only estrito. Esta dissertação investiga se criptogramas produzidos por Ascon-AEAD128, GIFT-COFB, Grain-128AEAD e Schwaemm256-128 podem ser distinguidos por classificadores de aprendizado de máquina quando o adversário observa apenas o texto cifrado. Para isso, foram gerados 180.000 criptogramas de 64~KB a partir de 300 chaves e 100 mensagens por chave, 80\% do Standardized Project Gutenberg Corpus e 20\% de imagens em tons de cinza do ImageNet-1k, cifradas de forma encadeada pelos quatro algoritmos e por dois controles, AES-ECB e um gerador de números pseudoaleatórios, com a mesma chave, o mesmo nonce e o mesmo texto em claro entre os seis, sob protocolo com particionamento por chave (240 chaves para treinamento e validação, 60 para teste) e seleção de atributos recalculada dentro de cada dobra da validação cruzada. Cinco representações do criptograma foram avaliadas: 641 descritores estatísticos, representação sequencial aprendida por CNN-1D, mapa de co-ocorrência de bigramas com CNN-2D, fusão híbrida e um Transformer hierárquico, submetidas aos classificadores Random Forest, SVM com núcleo RBF, LinearSVC, XGBoost e Regressão Logística, além de um meta-classificador por stacking sobre as saídas dos cinco caminhos. (ainda sem resultados). A dissertação disponibiliza, de todo modo, um protocolo e um dataset reprodutíveis de 180.000 amostras para trabalhos futuros de identificação algorítmica.
%fimmudança

 \textbf{Palavras-chave}: \imprimirpalavraschave
\end{resumo}

% resumo em inglês
\begin{resumo}[Abstract]
 \begin{otherlanguage*}{english}
%  \linespread{1.3}
\SingleSpacing
%iniciomudança
The standardization of the Ascon family by the National Institute of Standards and Technology (NIST), formalized in NIST SP 800-232, consolidated Lightweight Cryptography (LWC) as a response to the demand for authenticated encryption on devices with processing, memory, and energy constraints. However, the literature on cryptographic algorithm identification with machine learning reports high accuracies under protocols that allow information leakage between training and test, such as key reuse and feature selection outside the validation loop, and does not evaluate standardized or finalist lightweight AEAD schemes under a strict ciphertext-only scenario. This dissertation investigates whether ciphertexts produced by Ascon-AEAD128, GIFT-COFB, Grain-128AEAD, and Schwaemm256-128 can be distinguished by machine learning classifiers when the adversary observes only the ciphertext. For this, 180,000 ciphertexts of 64~KB were generated from 300 keys and 100 messages per key, 80\% from the Standardized Project Gutenberg Corpus and 20\% from grayscale ImageNet-1k images, encrypted in chained form by the four algorithms and by two controls, AES-ECB and a pseudorandom number generator, with the same key, nonce, and plaintext across the six, under a protocol with key-based partitioning (240 keys for training and validation, 60 for test) and feature selection recomputed inside each cross-validation fold. Five ciphertext representations were evaluated: 641 statistical descriptors, a sequential representation learned by a 1D CNN, a bigram co-occurrence map with a 2D CNN, a hybrid fusion, and a hierarchical Transformer, submitted to Random Forest, SVM with RBF kernel, LinearSVC, XGBoost, and Logistic Regression classifiers, plus a stacking meta-classifier over the outputs of the five paths. (no results yet). The dissertation makes available, in any case, a reproducible protocol and dataset of 180,000 samples for future work on algorithm identification.
%fimmudança
   \vspace{\onelineskip}

   \noindent
   \textbf{Keywords}: \imprimirkeywords
 \end{otherlanguage*}
\end{resumo}